\documentclass[10pt,a4paper]{amsart}
\usepackage[margin=19mm]{geometry}
\usepackage{amsmath,amssymb,mathtools}
\usepackage[scr=rsfs,cal=euler]{mathalfa}
\usepackage{xfrac}
\usepackage[unicode,hidelinks,bookmarks=false]{hyperref}
\newcommand{\ie}{i.e.\ }
\newcommand{\cf}{cf.\ }
\newcommand{\resp}{respectively\ }
\newcommand{\Subsection}{\subsection}
\providecommand{\nobreakdash}{\nobreak\hbox{-}}
\setlength{\emergencystretch}{2em}
\multlinegap0pt
\usepackage{bm}
\usepackage{bbm}
\usepackage{dsfont}
\usepackage{xspace}
\usepackage{cancel}
\usepackage{mathtools}
\usepackage{amsthm}
\makeatletter
\@ifundefined{theorem}{\newtheorem{theorem}{Theorem}}{}
\makeatother
\title[On the distribution of the telegraph meander and its propert]{On the distribution of the telegraph meander and its properties}
\author{Andrea Pedicone, Enzo Orsingher}
\date{Source: arXiv:2504.11387v2 (CC BY 4.0)}
\begin{document}
\maketitle

\noindent\emph{Source and licence.} On the distribution of the telegraph meander and its properties. Version arXiv:2504.11387v2 (CC BY 4.0), distributed under the Creative Commons Attribution 4.0 International licence, as recorded in the arXiv licence metadata for this exact version and shown on the abstract page https://arxiv.org/abs/2504.11387v2. Full rights notice: licenses/arxiv-2504.11387-CC-BY-4.0.txt.\par\smallskip

\noindent\textit{Editorial scope.} Counted unit: the Theorem whose source label is \texttt{pde telegraph meander} in the article (source0.tex lines 501--510, source label \texttt{pde telegraph meander}), delivered together with the article's own complete original proof of it (source0.tex lines 511--556, one \texttt{proof} environment of 46 lines). No mathematics is written, repaired or re-derived: statement and proof are the article's own text line for line. Required context retained uncounted: properties Bessel (lines 248--250); Telegraph meander t 4 (lines 268--270); properties derivative Bessel (lines 294--296). Every definition, display and label of those lines is retained unchanged. Omitted from this package and not redistributed: 1--247, 251--267, 271--293, 297--500, 557--817. Nothing inside a retained excerpt is omitted or altered: every retained body is delivered exactly as the article's lines are written, so no omission receipt of any kind accompanies this package. Citations: the retained excerpts carry no citation command at all, so no bibliography entry of the article is transcribed and no reference is redistributed. Reference adaptation: none was needed and none was made; every reference inside the delivered unit resolves inside it, so no unresolved reference or citation is produced. Printed slips kept literal: no printed word, symbol, hypothesis or number of the counted statement or of its proof was altered. Layout and class: the article's own class is \texttt{article} with packages including authblk, hyperref, amsmath, amsfonts, mathptmx, amsthm, geometry; the wrapper is the lane's pinned \texttt{amsart} editorial wrapper at 10pt on A4, which restates the theorem-like environments the retained text needs and adds the article's own macro definitions that the retained text needs (none), with extra wrapper packages (bm, bbm, dsfont, xspace, cancel, mathtools). The delivered numbering is the unit's own. Assets: no retained span contains a figure environment, a graphics-inclusion command, a PDF-page inclusion, a table or a TikZ picture; no image file is copied into this package, the package declares no asset dependency and no image file is redistributed with it. Licence: version arXiv:2504.11387v2 of this article is distributed under the Creative Commons Attribution 4.0 International licence, as recorded in the arXiv licence metadata for this exact version and shown on the abstract page https://arxiv.org/abs/2504.11387v2. A full rights notice is delivered at licenses/arxiv-2504.11387-CC-BY-4.0.txt. Characters: every retained excerpt is pure ASCII, so the delivered engine, XeLaTeX with its default Latin Modern text font, reports no missing character.
\begin{equation}\label{properties Bessel}
I_{\nu+1}(x) = I_{\nu-1}(x)-\frac{2\nu}{x}I_\nu(x).
\end{equation}

\begin{align}\label{Telegraph meander t 4}
&q_{\lambda,c}(x,t) = \lim_{\epsilon \downarrow 0}\mathsf{P}\{T(t) \in dx \lvert \min_{0\leq s \leq t} T(s) \geq -\epsilon,V(0) = -c\} =-\frac{\frac{\partial}{\partial x}p^{-}_{\lambda,c}(x,t)}{p_{\lambda,c}^{-}(0,t)} \;\; x \in [0,ct).
\end{align}

\begin{equation}\label{properties derivative Bessel}
\frac{d}{dx}I_\nu(x) = I_{\nu-1}(x) - \frac{\nu}{x}I_\nu(x).
\end{equation}

\begin{theorem}
Let $T^+ = (T^+(s))_{s \in [0,t]}$ be a telegraph meander. Then the absolutely continuous component of the law of $T^+(t)$ satisfies 
\begin{equation}\label{pde telegraph meander}
\frac{\partial^2u}{\partial t^2}-c^2\frac{\partial^2 u}{\partial x^2}+2\lambda\frac{\partial u}{\partial  t}-\frac{2a}{t}\frac{\partial u}{\partial t}=\Big(\frac{\lambda}{t}-\frac{2a}{ t^2}\Big)u
\end{equation}
where 
\begin{equation}
a(t) = \frac{I_1(\lambda t)}{I_0(\lambda t)+I_1(\lambda t)}.
\end{equation}
\end{theorem}

\begin{proof}
We rewrite equation \eqref{Telegraph meander t 4} in a suitable way
\begin{equation}\label{Telegraph meander t 5}
q_{\lambda,c}(x,t)p^{-}_{\lambda,c}(0,t) = -\frac{\partial}{\partial x}p^{-}_{\lambda,c}(x,t).
\end{equation}
Hence, by applying the second time derivative on both member of \eqref{Telegraph meander t 5} we obtain
\begin{align}\label{cp pde Telegraph meander 1}
&\frac{\partial^2}{\partial t^2}q_{\lambda,c}(x,t)p^{-}_{\lambda,c}(0,t)+2\frac{\partial}{\partial t}q_{\lambda,c}(x,t)\frac{\partial}{\partial t}p^{-}_{\lambda,c}(0,t)+q_{\lambda,c}(x,t)\frac{\partial^2}{\partial t^2}p^{-}_{\lambda,c}(0,t) =-\frac{\partial^3}{\partial t^2\partial x}p^{-}_{\lambda,c}(x,t)
\end{align}
Analogously, by taking in both members of \eqref{Telegraph meander t 5} the second order space derivative and also multiplying with $-c^2$, we get
\begin{align}\label{cp pde Telegraph meander 2}
-c^2\frac{\partial^2}{\partial x^2}q_{\lambda,c}(x,t)p^{-}_{\lambda,c}(0,t) = c^2\frac{\partial^3}{\partial x^3}p^{-}_{\lambda,c}(x,t).
\end{align}
We sum up equations \eqref{cp pde Telegraph meander 1} and \eqref{cp pde Telegraph meander 2} and it follows that
\begin{align}\label{cp pde Telegraph meander 3}
&\frac{\partial^2}{\partial t^2}q_{\lambda,c}(x,t)p^{-}_{\lambda,c}(0,t)+2\frac{\partial}{\partial t}q_{\lambda,c}(x,t)\frac{\partial}{\partial t}p^{-}_{\lambda,c}(0,t)+q_{\lambda,c}(x,t)\frac{\partial^2}{\partial t^2}p^{-}_{\lambda,c}(0,t) -c^2\frac{\partial^2}{\partial x^2}q_{\lambda,c}(x,t)p^{-}_{\lambda,c}(0,t)\nonumber\\&=-\frac{\partial}{\partial x}\Big\{\frac{\partial^2}{\partial t^2}p^{-}_{\lambda,c}(x,t)-\frac{\partial^2}{\partial x^2}c^2p^{-}_{\lambda,c}(x,t)\Big\}.
\end{align}
Now, the density function $p^{-}_{\lambda,c}(x,t)$ satisfies the telegraph equation
\begin{equation}\label{Telegraph equation}
\frac{\partial^2}{\partial t^2}p^{-}_{\lambda,c}(x,t) +2\lambda\frac{\partial}{\partial t}p^{-}_{\lambda,c}(x,t) =c^2\frac{\partial^2}{\partial x^2}p^{-}_{\lambda,c}(x,t),
\end{equation}
therefore, by means of \eqref{Telegraph meander t 5} and \eqref{Telegraph equation}, the right hand side of equation \eqref{cp pde Telegraph meander 3} becomes 
\begin{align*}
&2\lambda\frac{\partial^2}{\partial x\partial t}p^{-}_{\lambda,c}(x,t) = -2\lambda \frac{\partial}{\partial t}\Big\{q_{\lambda,c}(x,t)p^{-}_{\lambda,c}(0,t)\Big\}=-2\lambda\big\{\frac{\partial}{\partial t}q_{\lambda,c}(x,t)p^{-}_{\lambda,c}(0,t)+q_{\lambda,c}(x,t)\frac{\partial}{\partial t}p^{-}_{\lambda,c}(0,t)\big\}.
\end{align*}
As a result, equation \eqref{cp pde Telegraph meander 3} takes the form 
\begin{align}\label{cp pde Telegraph meander 5}
&\Big\{\frac{\partial^2}{\partial t^2}q_{\lambda,c}(x,t)-c^2\frac{\partial^2}{\partial x^2}q_{\lambda,c}(x,t)+2\lambda\frac{\partial}{\partial t}q_{\lambda,c}(x,t)\Big\}p^{-}_{\lambda,c}(0,t)+2\frac{\partial}{\partial t}q_{\lambda,c}(x,t)\frac{\partial}{\partial t}p^{-}_{\lambda,c}(0,t)\nonumber \\
&=-q_{\lambda,c}(x,t)\big\{\frac{\partial^2}{\partial t^2}p^{-}_{\lambda,c}(0,t)+2\lambda\frac{\partial}{\partial t}p^{-}_{\lambda,c}(0,t)\big\}.
\end{align}
We now compute the first and second time derivative of $p^{-}_{\lambda,c}(0,t)$. 
\begin{align}
&\frac{\partial}{\partial t}p^{-}_{\lambda,c}(0,t) = -\frac{\lambda^2}{2c}e^{-\lambda t}(I_0(\lambda t)+I_1(\lambda t))+\frac{\lambda^2}{2c}e^{-\lambda t}(I_0(\lambda t)+I_1(\lambda t)-\frac{1}{\lambda t}I_1(\lambda t)) \nonumber\\
&=-\frac{\lambda}{2ct}e^{-\lambda t}I_1(\lambda t) \label{First Derivative p-}\\
&\frac{\partial^2}{\partial t^2}p^{-}_{\lambda,c}(0,t) = \frac{\lambda}{2ct^2}e^{-\lambda t}I_1(\lambda t)+\frac{\lambda^2}{2ct}e^{-\lambda t}I_1(\lambda t)-\frac{\lambda^2}{2ct}e^{-\lambda t}(I_0(\lambda t)-\frac{1}{\lambda t}I_1(\lambda t)) \nonumber\\&=\frac{\lambda}{ct^2}e^{-\lambda t}I_1(\lambda t)+\frac{\lambda^2}{2ct}e^{-\lambda t}(I_1(\lambda t)-I_0(\lambda t)) \label{Second Derivative p-} ,\end{align}
where we made use of \eqref{properties Bessel} and \eqref{properties derivative Bessel}. Finally,
\begin{align}\label{cp Telegraph meander 6}
&\frac{\partial^2}{\partial t^2}p^{-}_{\lambda,c}(0,t)+2\lambda\frac{\partial}{\partial t}p^{-}_{\lambda,c}(0,t) = \frac{\lambda}{2c}e^{-\lambda t}\big\{2t^{-2}I_1(\lambda t)-\lambda t^{-1}\big(I_0(\lambda t)+I_1(\lambda t)\big)\big\}.
\end{align}
Therefore, by dividing with $p^{-}_{\lambda,c}(0,t)$ and substituting \eqref{First Derivative p-}, \eqref{Second Derivative p-} and \eqref{cp Telegraph meander 6} into \eqref{cp pde Telegraph meander 5}, we arrive at
\begin{align*}\label{pde Telegraph meander 2}
&\frac{\partial^2}{\partial t^2}q_{\lambda,c}(x,t) - c^2\frac{\partial^2}{\partial x^2}q_{\lambda,c}(x,t) +2\lambda\frac{\partial}{\partial t}q_{\lambda,c}(x,t)-\frac{2I_1(\lambda t)}{t(I_0(\lambda t)+I_1(\lambda t))} \frac{\partial}{\partial t}q_{\lambda,c}(x,t) \nonumber\\
&=-q_{\lambda,c}(x,t)\Big(\frac{2I_1(\lambda t)}{t^2(I_0(\lambda t)+I_1(\lambda t))}-\frac{\lambda}{t}\Big),
\end{align*}
that is \eqref{pde telegraph meander}.
\end{proof}

\end{document}
